\documentclass[11pt]{article}
\usepackage[a4paper,margin=25mm]{geometry}
\usepackage{amsmath,amssymb}
\usepackage[T1]{fontenc}
\usepackage{microtype}
\title{Optimized strategy for $[22][444][22]$}
\author{}
\date{}
\begin{document}
\maketitle
\section{Conventions and inequality}
\label{app:444_strategy}
The local Hilbert space is $\mathbb C^2\otimes\mathbb C^3\otimes\mathbb C^2$,
with tensor order $A,B,C$. Alice and Charlie each have two binary settings;
Bob has three settings with four outcomes each. All setting and outcome
indices start at zero. The state-vector basis is
\[
\mathcal E_{6a+2b+c}=|abc\rangle,\qquad
 a,c\in\{0,1\},\quad b\in\{0,1,2\}.
\]
Probabilities follow the Born rule,
\[
P(abc|xyz)=\operatorname{Tr}\!\left[\rho
 (A_{a|x}\otimes B_{b|y}\otimes C_{c|z})\right].
\]
The marginal notation is $P_{AB}(ab|xy)$ and $P_{BC}(bc|yz)$.
The inequality used for this strategy is
\begin{equation}
\begin{aligned}
S_{444}={}&2P_A(0|0)-P_B(0|0)+P_B(1|0)+P_B(2|1)\\
&-P_{AB}(02|00)-2P_{AB}(02|01)\\
&+P_{AB}(00|10)-P_{AB}(01|10)\\
&+P_{AB}(00|12)-P_{AB}(02|12)\\
&+P_{BC}(00|00)-P_{BC}(10|01)\\
&+P_{BC}(20|10)-P_{BC}(20|11)\\
&-P_{BC}(00|20)-P_{BC}(10|20)\\
&+P_{BC}(10|21)+P_{BC}(20|21)\geq -1.
\end{aligned}
\end{equation}
Its integer coefficient vector has $\ell_1$ norm $20$. We use
$S_{\rm norm}=S_{444}/20=c^Tp$, with $\|c\|_1=1$ and
$\beta_{\rm HI}=-1/20=-0.05$, independently checked by the HI support LP.
For the full-precision strategy below,
\[
\begin{aligned}
S_{\rm norm}&\simeq -0.050021278532119,\\
S_{444}&\simeq -1.000425570642376,\\
\Delta_{L_1}=\beta_{\rm HI}-S_{\rm norm}&\simeq 2.127853211881\times10^{-5}.
\end{aligned}
\]
This is an explicit locally optimized realization, not a certificate of a
global quantum optimum. It was obtained by continuing the previous
three-dimensional strategy in the full outer-loop search. The re-extracted
inequality agrees with the fixed inequality above within numerical precision.
All displayed entries are rounded to ten decimal places; full-precision
arrays should be used for numerical calculations.
\section{State}
The largest state amplitude is chosen real and positive, fixing only a
global phase. The density operator is $\rho=|\psi\rangle\langle\psi|$, where
\[
|\psi\rangle\simeq
\begin{pmatrix}
0.0051248927-0.0153029504\,i \\
0.6965009161 \\
-0.0100505039+0.0044208365\,i \\
-0.0006408368-0.0000800462\,i \\
0.0031300786+0.0003817264\,i \\
-0.0210317720+0.0029646895\,i \\
-0.0221644664+0.0039302707\,i \\
0.0010562264-0.0000022123\,i \\
-0.0151999594-0.0011606910\,i \\
-0.1664959713-0.6397100681\,i \\
-0.2709119781+0.0540788277\,i \\
0.0040452382+0.0044398641\,i
\end{pmatrix}.
\]
\section{Alice and Charlie}
\[
A_{0|0}\simeq \begin{pmatrix}
0.9998651065 & 0.0114257685+0.0020793111\,i \\
0.0114257685-0.0020793111\,i & 0.0001348902
\end{pmatrix}.
\]
\[
A_{1|0}=I_2-A_{0|0}.
\]
\[
A_{0|1}\simeq \begin{pmatrix}
0.4921194095 & 0.3690830832-0.3372173941\,i \\
0.3690830832+0.3372173941\,i & 0.5078805841
\end{pmatrix}.
\]
\[
A_{1|1}=I_2-A_{0|1}.
\]
\[
C_{0|0}\simeq \begin{pmatrix}
0.9978634076 & 0.0304106292-0.0347449480\,i \\
0.0304106292+0.0347449480\,i & 0.0021365839
\end{pmatrix}.
\]
\[
C_{1|0}=I_2-C_{0|0}.
\]
\[
C_{0|1}\simeq \begin{pmatrix}
0.0081812584 & 0.0441098255-0.0785407115\,i \\
0.0441098255+0.0785407115\,i & 0.9918187454
\end{pmatrix}.
\]
\[
C_{1|1}=I_2-C_{0|1}.
\]
\section{Bob}
To keep the complex matrices within the page width, define
\[
\mathcal H_3(d;z)=
\begin{pmatrix}
d_1&z_{12}&z_{13}\\
\overline z_{12}&d_2&z_{23}\\
\overline z_{13}&\overline z_{23}&d_3
\end{pmatrix}.
\]
Each list specifies one operator's diagonal and upper triangle; the lower
triangle follows from Hermiticity. Matrix indices refer to Bob's ordered
basis $(|0\rangle,|1\rangle,|2\rangle)$.
\subsection*{Setting $y=0$}
\[
\begin{gathered}
B_{0|0}\simeq\mathcal H_3(d;z),\\[3pt]
\begin{aligned}
d_1&=0.4262529703 \\
d_2&=0.4923299044 \\
d_3&=0.0814169482 \\
z_{12}&=0.3835025326-0.2505651425\,i \\
z_{13}&=-0.1861181072+0.0080159793\,i \\
z_{23}&=-0.1721637088-0.1021941566\,i
\end{aligned}
\end{gathered}
\]
\[
\begin{gathered}
B_{1|0}\simeq\mathcal H_3(d;z),\\[3pt]
\begin{aligned}
d_1&=0.4112640065 \\
d_2&=0.5076606803 \\
d_3&=0.0810751369 \\
z_{12}&=-0.3827310467+0.2495986983\,i \\
z_{13}&=-0.1824707676+0.0069061364\,i \\
z_{23}&=0.1740025659+0.1043156554\,i
\end{aligned}
\end{gathered}
\]
\[
\begin{gathered}
B_{2|0}\simeq\mathcal H_3(d;z),\\[3pt]
\begin{aligned}
d_1&=0.0812396510 \\
d_2&=0.0000047075 \\
d_3&=0.4187443655 \\
z_{12}&=-0.0003857341+0.0004832110\,i \\
z_{13}&=0.1842902163-0.0074608870\,i \\
z_{23}&=-0.0009194075-0.0010607251\,i
\end{aligned}
\end{gathered}
\]
\[
B_{3|0}=I_3-\sum_{b=0}^2 B_{b|0}.
\]
\subsection*{Setting $y=1$}
\[
\begin{gathered}
B_{0|1}\simeq\mathcal H_3(d;z),\\[3pt]
\begin{aligned}
d_1&=0.0005083636 \\
d_2&=0.3332886616 \\
d_3&=0.3328695682 \\
z_{12}&=0.0006597216-0.0037958809\,i \\
z_{13}&=0.0123530093+0.0013247756\,i \\
z_{23}&=-0.0000093785-0.0001435089\,i
\end{aligned}
\end{gathered}
\]
\[
\begin{gathered}
B_{1|1}\simeq\mathcal H_3(d;z),\\[3pt]
\begin{aligned}
d_1&=0.0005083636 \\
d_2&=0.3332886616 \\
d_3&=0.3328695682 \\
z_{12}&=0.0006597216-0.0037958809\,i \\
z_{13}&=0.0123530093+0.0013247756\,i \\
z_{23}&=-0.0000093785-0.0001435089\,i
\end{aligned}
\end{gathered}
\]
\[
\begin{gathered}
B_{2|1}\simeq\mathcal H_3(d;z),\\[3pt]
\begin{aligned}
d_1&=0.9984749093 \\
d_2&=0.0001340125 \\
d_3&=0.0013912942 \\
z_{12}&=-0.0019791649+0.0113876424\,i \\
z_{13}&=-0.0370590278-0.0039743268\,i \\
z_{23}&=0.0000281308+0.0004305370\,i
\end{aligned}
\end{gathered}
\]
\[
B_{3|1}=I_3-\sum_{b=0}^2 B_{b|1}.
\]
\subsection*{Setting $y=2$}
\[
\begin{gathered}
B_{0|2}\simeq\mathcal H_3(d;z),\\[3pt]
\begin{aligned}
d_1&=0.4588197756 \\
d_2&=0.4922179943 \\
d_3&=0.0000509373 \\
z_{12}&=0.3984001506-0.2590688479\,i \\
z_{13}&=0.0025547882-0.0041027502\,i \\
z_{23}&=0.0045349467-0.0021199393\,i
\end{aligned}
\end{gathered}
\]
\[
\begin{gathered}
B_{1|2}\simeq\mathcal H_3(d;z),\\[3pt]
\begin{aligned}
d_1&=0.0002034409 \\
d_2&=0.0000061535 \\
d_3&=0.9964162677 \\
z_{12}&=-0.0000057475-0.0000349062\,i \\
z_{13}&=0.0132497951+0.0052108435\,i \\
z_{23}&=-0.0012684119+0.0021261937\,i
\end{aligned}
\end{gathered}
\]
\[
\begin{gathered}
B_{2|2}\simeq\mathcal H_3(d;z),\\[3pt]
\begin{aligned}
d_1&=0.4429781571 \\
d_2&=0.5077580187 \\
d_3&=0.0000220699 \\
z_{12}&=-0.3980503291+0.2578403483\,i \\
z_{13}&=0.0027396651-0.0015031660\,i \\
z_{23}&=-0.0033367378-0.0002439412\,i
\end{aligned}
\end{gathered}
\]
\[
B_{3|2}=I_3-\sum_{b=0}^2 B_{b|2}.
\]
\section{Numerical checks}
The full-precision maximum POVM residual is $1.110\times10^{-15}$,
and the minimum effect eigenvalue is $-8.523\times10^{-18}$
(zero within floating-point accuracy). The rounded state is renormalized
when checking the displayed values, and final outcomes are reconstructed
by completeness. This gives
\[
\Delta_{L_1}^{\rm displayed}\simeq 2.127853295427\times10^{-5}.
\]
Rounding may create tiny negative effect eigenvalues. The accompanying
full-precision numerical strategy is the reference for positivity checks.
\end{document}
